\documentclass[a4paper,10pt]{article}
\usepackage[top=1cm, bottom=1cm, left=1.2cm, right=1.2cm]{geometry}
\usepackage{array}
\usepackage[dvipsnames,table]{xcolor}
\usepackage[hidelinks]{hyperref}
\usepackage{enumitem}
\usepackage{titlesec}
\usepackage{paracol}
\usepackage{microtype}
\usepackage[T1]{fontenc}
\usepackage[utf8]{inputenc}
\usepackage{helvet}
\renewcommand{\familydefault}{\sfdefault}
\usepackage{graphicx}
\usepackage{tikz}

% =============================================
% ACCENT COLOR – change RGB values to your liking
% Examples: blue (0,90,160) | green (40,120,70) | gray (60,60,60)
% =============================================
\definecolor{accentcolor}{RGB}{120, 47, 64}

\pagestyle{empty}
\setlength{\parindent}{0pt}
\setlength{\columnsep}{0.8cm}

\titleformat{\section}
  {\normalfont\normalsize\bfseries}
  {}{0em}{\color{accentcolor}}
  [\vspace{-2pt}\color{accentcolor}\rule{\linewidth}{0.6pt}\vspace{2pt}]
\titlespacing{\section}{0pt}{6pt}{4pt}

\newcommand{\cvitem}[1]{%
  \begin{itemize}[leftmargin=1em, itemsep=0pt, topsep=1pt, parsep=0pt]
    \small #1
  \end{itemize}
}

% =============================================
% LOGO ENTRY
% Usage: \logoentry{logo_file.png}{Title | City}{Subtitle | Dates}
% To resize all logos: change height=1.3cm below
% =============================================
\newcommand{\logoentry}[3]{%
  \begin{minipage}[c]{0.18\linewidth}
    \includegraphics[width=\linewidth,height=1.3cm,keepaspectratio]{#1}
  \end{minipage}%
  \hspace{5pt}%
  \begin{minipage}[c]{0.78\linewidth}
    {\small\bfseries #2}\\
    {\footnotesize #3}
  \end{minipage}%
}

\begin{document}

% =============================================
% HEADER
% Replace with your own info
% =============================================
\begin{minipage}[c]{0.1\textwidth}
  \includegraphics[width=\linewidth]{qr-code.png}
  % Generate your LinkedIn QR code and save as qr-code.png
\end{minipage}
\hfill
\begin{minipage}[c]{0.73\textwidth}
  \begin{center}
    {\LARGE\bfseries Firstname LASTNAME}\\[3pt]
    {\small\itshape Your job title and a short intro sentence. Mention your experience level, main expertise, and what you are looking for.}\\[5pt]
    \small
    +XX X XX XX XX XX \quad
    your.email@example.com \quad
    \href{https://github.com/yourusername}{github.com/yourusername}
  \end{center}
\end{minipage}
\hfill
\begin{minipage}[c]{0.12\textwidth}
  \begin{tikzpicture}
    \clip (0,0) circle (1.3cm);
    \node at (0,0.3) {\includegraphics[width=3cm]{photo.jpg}};
    % Replace photo.jpg with your own square photo
  \end{tikzpicture}
\end{minipage}

\vspace{2pt}
{\color{accentcolor}\rule{\linewidth}{1pt}}
\vspace{3pt}

\setlength{\columnseprule}{0pt}
\begin{paracol}{2}

%========== LEFT COLUMN ==========%

\section{Education}

% --- Degree 1 ---
\logoentry{school1_logo.png}{University or School Name | City}{Degree Title, Specialization | YYYY – YYYY}
\cvitem{
  \item Relevant coursework, specialization topics, or technologies studied.
  \item \textbf{Key projects:} Project Name 1, Project Name 2, Project Name 3.
}

\vspace{4pt}
% --- Exchange semester (optional) ---
\logoentry{school2_logo.png}{Exchange University | City, Country}{Exchange Semester, Level | Month YYYY – Month YYYY \\ \textit{Notable ranking or distinction}}
\cvitem{
  \item Courses taken: Course 1, Course 2, Course 3, Language studied.
}

\vspace{4pt}
% --- Degree 2 ---
\logoentry{school3_logo.png}{University Name | City}{Bachelor's Degree in Field – Honours | YYYY – YYYY \\ \textit{Optional note, e.g. double major or specialization}}
\cvitem{
  \item \textbf{Key projects:} Project Name 1, Project Name 2.
}

\vspace{3pt}
\section{Technical Skills}

% Add or remove categories as needed. Bold your strongest skills.
{\small
\textbf{Programming:} \textbf{Language1}, \textbf{Language2}, Language3, Language4, Language5.\\[2pt]
\textbf{Web \& Frontend:} \textbf{HTML}, \textbf{CSS}, Framework1, Framework2.\\[2pt]
\textbf{Data \& ML:} Tool1, Tool2, Tool3.\\[2pt]
\textbf{Domain:} Skill1, Skill2, \textbf{Skill3}.\\[2pt]
\textbf{Other:} Git, Linux, Tool4.
}

\vspace{3pt}
\section{Associations \& Hackathons}
\cvitem{
  \item Association Name – Role (e.g. Active member, President)
  \item Hackathon Name – Year
  \item Hackathon Name – Year
}

\vspace{3pt}
\section{Soft Skills}
{\small Skill1 · Skill2 · Skill3 · Skill4 · Skill5}

\vspace{3pt}
\section{Languages}
{\small
\textbf{Language1} – Native \quad \textbf{Language2} – Level/Score \quad \textbf{Language3} – Level
}

%========== RIGHT COLUMN ==========%
\switchcolumn

\section{Professional Experience}

% --- Experience 1 (most recent) ---
\logoentry{company1_logo.png}{Company or Institution Name | City, Country}{Job Title | Month YYYY – Month YYYY}

\vspace{2pt}
{\small Brief description of the role and main achievements. Keep it to 2–3 sentences, focusing on impact and technologies used.}

\vspace{5pt}
% --- Experience 2 ---
\logoentry{company2_logo.png}{Company Name | City}{Job Title | Month YYYY – Month YYYY}

\vspace{1pt}
{\small\textbf{Category 1 (e.g. Risk Management):}}
\cvitem{
  \item Responsibility or achievement. Use \textbf{bold} to highlight key terms.
  \item Responsibility or achievement involving a specific tool or methodology.
  \item Quantified result if possible (e.g. reduced X by Y\%).
  \item Cross-functional or international collaboration example.
}

{\small\textbf{Category 2 (e.g. Awareness \& Training):}}
\cvitem{
  \item Responsibility or achievement. Highlight \textbf{key actions}.
  \item Responsibility or achievement.
  \item Responsibility or achievement.
}

{\small\textbf{Category 3 (e.g. Coordination \& Governance):}}
\cvitem{
  \item Responsibility involving frameworks or standards (e.g. \textbf{ISO 27001}).
  \item Tool or process you maintained or improved.
  \item Stakeholder or cross-team coordination example.
}

\vspace{4pt}
% --- Experience 3 ---
\logoentry{company3_logo.png}{Company Name | City}{Internship Title | Month YYYY – Month YYYY}

\vspace{2pt}
{\small Brief description of internship tasks and skills applied. Mention technologies or tools used.}

\vspace{3pt}
\section{Interests}
\cvitem{
  \item Interest or hobby 1 (with any notable achievement).
  \item Interest or hobby 2 (with any notable achievement).
  \item Interest or hobby 3.
  \item Interest or hobby 4.
}

\end{paracol}
\end{document}
